% !TeX spellcheck = en_US
% 特殊注释, 告知编译器使用美国英语作为拼写检查的语言

\documentclass[UTF8,a4paper,twoside]{ctexrep}  % 选项分别是UTF-8编码, A4纸张大小, 双面打印
% \documentclass[a4paper,UTF8,fontset=none]{ctexrep}  % fontset=none是为了后面指定自定义字体

\usepackage{amsmath,amssymb,amsfonts}  % 数学相关
\usepackage{ctex}
\usepackage{geometry}
\usepackage{makeidx}
\usepackage{graphicx}
\usepackage{caption}
\usepackage{verbatim}

\usepackage[colorlinks=true,allcolors=black]{hyperref}  % 链接相关, 设置为有颜色, 颜色是黑色, 下面选项类似
%\usepackage[colorlinks=true,linkcolor=black,citecolor=black，urlcolor=black]{hyperref}
%\usepackage[colorlinks=true]{hyperref}

\usepackage[nameinlink]{cleveref}  % 针对章节、图表和公式的**标签**
%\usepackage[labelsep=space,labelformat=dash]{caption}  % 用于设置图表和表格标题的格式, 分别是标签和文本有空格和标签和文本有破折号

\makeindex

% ctex
% \setCJKfamilyfont{BoldFont=SimHei,ItalicFont=KaiTi}{SimSun}  % 粗体使用SimHei, 斜体使用KaiTi
\setCJKmainfont[BoldFont=SimHei]{SimSun}  % 使用SimHei作为中文粗体  % 对应rm罗马
\setCJKsansfont{SimHei}  % 使用SimHei作为中文无衬线字体  % 对应sf无衬线
\setCJKmonofont{FangSong}  % 使用FangSong作为中文等宽字体  % 对应tt打字机
% fontspec, 是XeLaTeX和LuaLaTeX字体选择工具
\setmainfont{Times New Roman}  % 使用Times New Roman作为主要西文字体
\setmonofont{Times New Roman}  % 使用Times New Roman作为等宽西文字体


\geometry{left=3.18cm,right=3.18cm,top=2.54cm,bottom=2.54cm}  % 设置页边距
\linespread{1.5}  % 设置行距 (latex基础行距默认为1.2, 而word中默认为1.3)

\ctexset{
	chapter={
		name={第,章},
		number=\arabic{chapter},
		format=\raggedright\bfseries\zihao{3},
		beforeskip=0ex,
		afterskip=0.5ex,
		aftername=\enspace,
	},
	section={
		format=\raggedright\bfseries\zihao{4},
		beforeskip=0.5ex,
		afterskip=0ex,
		aftername=\enspace,
	},
	subsection={
		format=\raggedright\bfseries\zihao{-4},
		beforeskip=1ex,
		afterskip=1ex,
		aftername=\enspace,
	},
	contentsname={\hfill\zihao{3}\textbf{目 \quad 录}\hfill}
}

% \usepackage{titlesec}  % 设置标题样式 
% \titleformat{\chapter}{\zihao{3}\bfseries}{第\,\arabic{chapter}\,章}{\enspace}{}  % 设置一级标题左对齐，中文字体为黑体，英文字体为Times New Roman，字号为三号

% \DeclareCaptionLabelFormat{dash}{#1 \arabic{chapter}--\arabic{figure}}
% \DeclareCaptionLabelSeparator{dash}{-}
\renewcommand{\figurename}{图}
\renewcommand{\tablename}{表}
\crefname{figure}{图}{图}
\crefname{table}{表}{表}
\renewcommand{\thefigure}{\arabic{chapter}--\arabic{figure}}
\renewcommand{\thetable}{\arabic{chapter}--\arabic{table}}
% \makeatletter
% \crefformat{figure}{#2\cref@figure@name~\thefigure#3}
% \makeatother
% \captionsetup[figure]{labelsep=space,labelformat=dash,font={small,bf}}
\DeclareCaptionFont{five}{\zihao{5}\bfseries}
\captionsetup[figure]{labelsep=space,font=five}
\captionsetup[table]{labelsep=space,font=five,position=top}

\usepackage{titling}
% Title Page
\title{论文题目}
%\author{作者姓名}

\usepackage{fancyhdr}  % 定义页眉页脚

\fancypagestyle{tocstyle}{
	\fancyhf{} % clear all header and footer fields
	\renewcommand{\headrulewidth}{0pt}
	\fancyfoot[C]{\zihao{-5}\thepage}
}
\fancypagestyle{mainstyle}{
	\fancyhf{}
	\renewcommand{\headrulewidth}{0.3pt}
	\fancyhead[CO]{\zihao{5}\texttt{浙江理工大学本科毕业设计（论文）}}
	\fancyhead[CE]{\zihao{5}\texttt{\thetitle}}
	\fancyfoot[CE,CO]{\zihao{-5}\thepage}
}

\usepackage{placeins}
\usepackage{xcolor}
\usepackage[cache=false]{minted}
\setminted{
	frame=lines,
	framesep=2mm,
	baselinestretch=1.2,
%	bgcolor=lightgray,
	fontsize=\footnotesize,
	linenos
}

\graphicspath{{figure/}}

\usepackage{tocloft}
\setlength\cftbeforetoctitleskip{0ex}
\setlength\cftaftertoctitleskip{0ex}
\setlength\cftbeforechapskip{0ex}
\setlength\cftbeforepartskip{0ex}
\renewcommand{\cftpartfont}{\zihao{4}\bfseries}
\renewcommand{\cftsecfont}{\zihao{5}\mdseries}
\renewcommand{\cftsubsecfont}{\zihao{5}\ttfamily}
\renewcommand{\cftchappagefont}{\mdseries}
\renewcommand{\cftchapleader}{\cftdotfill{\cftdotsep}}

\usepackage{pdfpages}

% GB/T 7714-2015标准
\usepackage[backend=biber,style=gb7714-2015,gbnoauthor=true]{biblatex}
\addbibresource{thesis.bib}

\begin{document}
\zihao{-4}

\includepdf[pages=-]{诚信声明.pdf}
\phantomsection
\pdfbookmark[0]{摘要}{bookmark:摘要}
\cftaddtitleline{toc}{part}{摘 \quad 要}{}
\renewcommand{\abstractname}{\textbf{\zihao{-3}摘 \quad 要}}
\begin{abstract}

    这里是摘要。

    可以有多个自然段。

    一个自然段足够长会是什么效果呢呢呢呢呢呢呢呢呢呢呢呢呢呢呢呢呢呢呢呢呢呢呢呢呢呢呢呢呢呢呢呢呢呢呢呢呢呢呢呢呢呢呢呢呢呢呢呢呢呢？
	
	\vspace{0.5ex}
	\noindent\textbf{关键词：}关键词1；关键词2；关键词3；关键词4；关键词5
\end{abstract}
\clearpage


\phantomsection
\pdfbookmark[0]{Abstract}{bookmark:Abstract}
\cftaddtitleline{toc}{part}{Abstract}{}
\renewcommand{\abstractname}{\textbf{\zihao{-3}Abstract}}
\begin{abstract}
    % 这里的\qquad必须有, 否则英文摘要第一自然段不能正确缩进
	\qquad Here is the abstract.
    
    % 第二段则不需要
    There can be multiple natural segments.
    
    What is the effect of a natural segment being long enoughhhhhhhhhhhhhhhhhhhhhhhhhhhhhhhhhhhhhhhhhhhhhhhhhh?

	\vspace{0.5ex}
	\noindent\textbf{Keywords: }Keywords1; Keywords2; Keywords3; Keywords4; Keywords5
\end{abstract}
\clearpage


\pdfbookmark[0]{目录}{bookmark:目录}
\pagestyle{tocstyle}
\fancypagestyle{plain}{
    \pagestyle{tocstyle}
}
\pagenumbering{Roman}
\tableofcontents

\clearpage
\pagestyle{mainstyle}
\fancypagestyle{plain}{
    \pagestyle{mainstyle}
}
\pagenumbering{arabic}

% \include{chapter/introduction}
% 
% {
\ctexset{chapter/format=\centering\bfseries\zihao{3}}
\chapter*{参考文献}

\phantomsection
\pdfbookmark[0]{参考文献}{bookmark:参考文献}
\cftaddtitleline{toc}{chapter}{参考文献}{\thepage}
\printbibliography[heading=none]
\clearpage
}
% {
\ctexset{chapter/format=\centering\bfseries\zihao{3}}
\chapter*{致 \quad 谢}

\phantomsection
\pdfbookmark[0]{致谢}{bookmark:致谢}
\cftaddtitleline{toc}{chapter}{致 \quad 谢}{\thepage}

致谢内容

\begin{flushright}
姓名

xxxx 年 xx 月 xx 日
\end{flushright}
}

\end{document}
